
% !TEX program = lualatex
\documentclass[a4paper,12pt]{bxjsbook}

% =========================================================
% Packages
% =========================================================

% --- Japanese (LuaLaTeX) ---
\usepackage{luatexja}

% --- math / layout ---
\usepackage{amssymb}
\usepackage{amsthm}
\usepackage{mathtools}
\usepackage{geometry}
\usepackage{mathpartir}
\geometry{margin=25mm}

% --- lists / links / figures ---
\usepackage{enumitem}
\usepackage{xurl}
\usepackage{hyperref}
\usepackage{tikz}
\usetikzlibrary{arrows.meta, positioning, calc, fit, backgrounds, decorations.markings}

\hypersetup{
  colorlinks=true,
  linkcolor=blue,
  urlcolor=blue,
  citecolor=blue
}

% =========================================================
% length definition
% =========================================================
\newlength{\customindent}
\setlength{\customindent}{2em}

% =========================================================
% macro definition
% =========================================================
\newcommand{\BodyListFixBegin}{%
  \begingroup
  \setlength{\rightskip}{0pt}%
  \setlength{\parindent}{0pt}%
  % \setlength{\displaywidth}{\linewidth}%
}
\newcommand{\BodyListFixEnd}{\ifhmode\par\fi\endgroup}

\newenvironment{BodyListFix}{\BodyListFixBegin}{\BodyListFixEnd}

% =========================================================
% environment definition
% =========================================================

% Main_Sentense styles
\newenvironment{body}{
  \par
  \begingroup
  \setlength{\leftskip}{\customindent}
  \setlength{\rightskip}{0pt}
  \setlength{\parindent}{1em}
  \setlength{\displayindent}{\customindent}
  \setlength{\displaywidth}{\dimexpr\linewidth-\customindent\relax}
}{
  \par
  \endgroup
}

% description in Main_Sentense styles
\newlist{bodydescription}{description}{1}
\setlist[bodydescription]{
  labelwidth=0pt,
  labelsep=0pt,
  itemindent=0pt,
  style=nextline,
  font=\bfseries
}

% itemize in Main_Sentense styles
\newlist{bodyitemize}{itemize}{3}
\setlist[bodyitemize]{
  leftmargin=3em,
  label=\textbullet,
  before=\begin{BodyListFix},
  after=\end{BodyListFix}
}
\setlist[bodyitemize,2]{leftmargin=3.5em}
\setlist[bodyitemize,3]{leftmargin=4em}

% enumerate in Main_Sentense styles
\newlist{bodyenumerate}{enumerate}{3}
\setlist[bodyenumerate]{
  leftmargin=4em,
  label=(\arabic*),
  before=\begin{BodyListFix},
  after=\end{BodyListFix}
}
\setlist[bodyenumerate,2]{leftmargin=4.5em, label=(\alph*)}
\setlist[bodyenumerate,3]{leftmargin=5em, label=(\roman*)}

% Theorem styles
\newtheoremstyle{theorem}
  {\baselineskip}
  {\baselineskip}
  {
    \normalfont
    \setlength{\leftskip}{\customindent}
    \setlength{\parindent}{0pt}
    \setlength{\rightskip}{0pt}
    \setlength{\displayindent}{\customindent}
    \setlength{\displaywidth}{\dimexpr\linewidth-\customindent\relax}
  }
  {0pt}
  {\bfseries}
  {.}
  {\newline}
  {
    \hskip-\customindent
    \thmname{#1}\thmnumber{ #2}\thmnote{（#3）}
  }

\theoremstyle{theorem}
\newtheorem{theorem}{定理}[chapter]
\newtheorem{lemma}[theorem]{補題}
\newtheorem{proposition}[theorem]{命題}
\newtheorem{corollary}[theorem]{系}
\newtheorem{definition}[theorem]{定義}
\newtheorem{example}[theorem]{例}
\newtheorem{remark}[theorem]{注意}
\newtheorem{abbreviation}[theorem]{略記}

% description in Theorem styles
\newlist{theoremdescription}{description}{3}
\setlist[theoremdescription]{
  leftmargin=5em,
  label=\textbullet,
  before=\begin{BodyListFix},
  after=\end{BodyListFix}
}
\setlist[theoremdescription,2]{leftmargin=5.5em}
\setlist[theoremdescription,3]{leftmargin=6em}

% itemize in Theorem styles
\newlist{theoremitemize}{itemize}{3}
\setlist[theoremitemize]{
  leftmargin=5em,
  label=\textbullet,
  before=\begin{BodyListFix},
  after=\end{BodyListFix}
}
\setlist[theoremitemize,2]{leftmargin=5.5em}
\setlist[theoremitemize,3]{leftmargin=6em}

% enumerate in Theorem styles
\newlist{theoremenumerate}{enumerate}{3}
\setlist[theoremenumerate]{
  leftmargin=6em,
  label=(\arabic*),
  before=\begin{BodyListFix},
  after=\end{BodyListFix}
}
\setlist[theoremenumerate,2]{leftmargin=6.5em, label=(\alph*)}
\setlist[theoremenumerate,3]{leftmargin=7em, label=(\roman*)}

% =========================================================
% Proof Environment Setting
% =========================================================
\makeatletter
\renewenvironment{proof}[1][\proofname]{%
  \par\pushQED{\qed}%
  \normalfont
  \topsep6\p@\@plus6\p@\relax
  \list{}{%
    \leftmargin=\customindent
    \rightmargin=0pt
    \labelwidth=0pt
    \labelsep=0pt
    \itemindent=0pt
    \listparindent=0pt
    \parsep=0pt
  }%
  \item[\itshape #1\@addpunct{.}\hspace{0.5em}]%
  \setlength{\displayindent}{\customindent}%
  \setlength{\displaywidth}{\dimexpr\linewidth-\customindent\relax}%
  \ignorespaces
}{%
  \popQED\endlist\@endpefalse
}
\makeatother

% =========================================================
% Convenience commands
% =========================================================
\newcommand{\addchaptertotoc}[1]{%
  \chapter*{#1}%
  \addcontentsline{toc}{chapter}{#1}%
}
\newcommand{\dotminus}{\mathbin{\dot{-}}}

\DeclareMathOperator{\AffHull}{AffHull}
\DeclareMathOperator{\AffComb}{AffComb}


% =========================================================
% Metadata
% =========================================================
\title{チェバの定理の一般化と\\ Lean4/mathlib4への形式化}
\author{
秋田 隼\\
早稲田大学\\
基幹理工学部 数学科B4\\
1W221002-0
}

\date{$2026$/$1$/$9$}

% =========================================================
% Document
% =========================================================
\begin{document}

    \maketitle

    % \addchaptertotoc{進捗(最終的に消します)}
    % \textbf{必要十分条件を主張する定理として一般化を行いたかったが、どのバージョンの必要十分形が最も汎用的で理想のチェバの定理と言えるのか結論が出ず、また仮に結論が出たとして、おそらく実装が間に合わないことから、ジョセフマイヤーズ氏の定理をそのまま逆にした定理を示し実装することにしました。もし余裕があればさらに一般化を行いたいと思います。}
    % 2026/1/22 2.3 Lean4・mathlib4の説明を編集
    % 2026/1/30 一般化とは何か、古典的チェバの主張になぜパターンがあるかの説明を補足、定義の整合性や証明部分はまだ未完成です

    \frontmatter
    \addchaptertotoc{要旨}

    \begin{body}
        \indent 本研究は、古典的チェバの定理を出発点として、チェバの定理の一般化を定式化し、関連する補題のLean4/mathlib4上での形式化を目標とする。

        \vskip\baselineskip

        \indent 主な貢献は以下である。

        \begin{bodyitemize}[leftmargin=3em]
            \item チェバの定理の一般化の定式化の紹介。
            \item 既存のmathlib4の構造を調査し、再利用可能部分と不足部分を切り分ける設計指針。
            \item 不足する補題・定義をモジュール化して追加し、定理の機械検証を通す。
        \end{bodyitemize}

        \vskip\baselineskip
        \indent 実装はGitHubリポジトリで公開している：
        \par\noindent
        \url{https://github.com/Aj1905/LeanResearch.git}
    \end{body}

    \tableofcontents

    \mainmatter

    % =========================================================
    \chapter{序論}
    % =========================================================

    \section{研究背景と動機}
    \begin{body}
        \indent 本研究は、mathlib4に存在する順方向チェバの定理の形式化を基盤として、その定理の逆をLean4上で構成し、mathlib4への追加を目指す。mathlib4は、Lean Prover Communityにより開発される大規模数学ライブラリ\cite{LeanProverCommunity}であり、分野により整備状況には濃淡がある。
        特に幾何学は、形式化が代数的手法に比べて技術的に困難であることや、図形的直観の形式言語翻訳の複雑さが起因するためか発展が停滞気味である。こうした開発目標を可視化すべく、数学において重要でありながら未だ形式化されていない定理が「missing 100 theorems」としてリストアップされている。

        \indent チェバの定理は、図形の共点性を代数条件へ翻訳する基本的な道具であり、多数の幾何学的議論の分岐点として機能する。この種の定理を一般化してライブラリ化することの価値は、定理本体に加えて単体・面・アフィン結合・重心座標などの重要な周辺概念を再利用可能な形で整備できる点にある。よって今後の幾何定理を形式化する際の基盤を構築できるという点で大きな意義を持つと考える。

        \indent さらに、形式化は教科書的な証明に潜む、ゼロ除算回避、図形が退化する場合の排除、必要な代数構造などの暗黙の仮定を明示し、どの仮定のもとでどの結論が成立するかを定理の型として確定する。このことは、一般化の汎用性を証明の文章ではなくLeanの型と依存関係によって機械的に検証可能な形で提示できる点で、数学的にも工学的にも意義がある。近年もチェバ型定理の一般化は研究されており、本研究はその形式的基盤をmathlib4に提供する試みとして位置づけられる。
    \end{body}

    \section{本稿の構成}
    \begin{body}
        \indent 第$2$章でLean4について解説し、第$3$章で一般化とは何か基礎論的な視点から眺める。第$4$章でチェバの定理とアフィン幾何の関連を議論し、第$5$章で古典的チェバの定理およびその逆がどのように一般化され得るかを考える。第$6$章で実際に行った実装の詳細を見る。
        \indent 本稿の議論は暗黙に、自然演繹を推論規則として採用した古典一階述語論理体系を用いて行なっている。集合論に言及する時、その概念の定義はKunenに従うものとする。\cite{Kunen}
    \end{body}


    % =========================================================
    \chapter{Lean4による形式化}
    % =========================================================
    \section{形式言語の概要}
    \begin{body}
        \indent 自然言語の主張はさまざまな点で曖昧さを含む。定義が省略されていたり、暗黙の前提があったり、同じ記号が別の意味を持っていたりすることで本来したいはずの主張が伝わらないことや別の意味で解釈されてしまうこともある。

        \indent これを解決するために用いられるのが\emph{形式言語}と呼ばれる定義が厳密に定まった言語体系であり、自然言語の文を形式言語に翻訳する作業は\emph{形式化}と呼ばれる。一度形式化された文章は、読み手がその体系の定義を正確に把握している限り、いつ誰がどこで読んでも意味が一意に定まる。これはコンピュータが文章を読む場合も例外ではない。この形式言語の文章をどう解釈しどう推論するかも含めて厳密に定義したものを形式体系と呼ぶ。

        \indent 一方数学とは、過去の数学者の成果の積み重ねからなる学問である。数学的主張は、他者に共有され、検証され、再利用される形で後世に伝え残されてはじめて真価を発揮する。このとき、時代や場所、読み手によって意味や解釈が揺れることは望ましくない。いついかなる時も同じ主張をすることが求められる。まさに形式言語の特徴が最大限活かされる分野の一つである。現代数学は集合論を一階述語論理で公理化した形式体系（主にZFやZFC）を基礎として用いて構築されことが多い。

        \indent しかし一度考え直してみると、数学を記述する形式言語は必ずしも一階述語論理に限る必要はない。一階述語論理は確かに人間が数学を理解する上で相性のいい形式言語の一種かもしれない。形式言語の中では比較的読みやすく、また歴史的に標準化されているため慣れてもいる。しかし本来、言語は読み手と目的に最適化することで意味伝達という言語本来の役割を最大限発揮する。伝えたい意味内容に対し、読み手と目的を明確にして最適な形式言語を選択することが重要となる。そのための一階述語論理の代替案の一つとして挙げられるのが型理論である。

        \indent 型理論とは、命題を型、証明をその型の項として表現するという特徴を持った形式体系である。集合論において全ての数学的対象が集合であるように、型理論において全ての数学的対象は型を持つ。詳細な説明は\hyperref[appendix:type-theory]{付録}にまとめた。たとえば命題 $P$ の証明は $p:P$（pは型Pの項）という形で与えられる。項$p$を構築する厳密な規則が用意されているため、その規則に従っているかを機械的に確認する作業を通して証明の正しさを検証できる。その意味で、コンピュータにとって検証のしやすい言語であると言える。

    \end{body}

    \section{証明支援系の概要}
    \begin{body}
        \indent 証明支援系とは、数学の定義・定理・証明を形式言語で記述することに特化したプログラミング言語、かつその言語で記述されたコードが公理や推論規則に照らして正しいかを機械的に検証する機能を備えたもの、の総称である。
        
        \indent 「証明を自動で発見する」ことを目的としたATP（Automated Theorem Proving）と、「与えられた証明が正しいことを厳密に検査する」ことを目的としたITP（Interactive Theorem Proving）が存在する。

        \indent 証明支援系には基礎となっている形式体系の種類に応じて複数の系統がある。代表的なものとして、依存型理論に基づくLeanやCoq（現在Rocqに改名）、高階論理系のIsabelle、HOLなどが存在し、目的に応じて使い分けられている。

        \indent この技術は応用面では、まずソフトウェアの形式検証に用いられる。ソフトウェアは各々が仕様という設計目標を持つ。この仕様を数学的な命題として書き下すことで、実装されたプログラムがその仕様を満たしているか否かという問題を、数学的な真偽判定の問題に帰着させることができる。挙動をあらかじめ検証しておくことで、信頼性の高い堅牢なソフトウェア開発が可能となる。

        \indent また近年、証明支援系はLLMとも結びつき始めている。LLMはIMO 2025(国際数学オリンピック)において金メダル基準相当の成績に達したとする報告も出ており\cite{DeepMindIMO2025}、自然言語推論の性能が飛躍的に向上している。一方で、もっともらしいが誤りを含む出力（ハルシネーション）は依然としてLLMの課題である。
        しかし、近年注目されている定理証明器を統合した枠組みでは、自然言語の解答案全体または一部を形式化して検証し、失敗時にはフィードバックにより解答を再生成するというループが可能となる。そのため、形式検証が及ぶ範囲についてはハルシネーションを大幅に抑制でき、結果として出力全体に対するハルシネーションの頻度も抑えられる。これは、ハルシネーション自体を抑えるのではなく正確な解答を返すまで出力を繰り返させるという発想の転換である
        （下図参照、問題文の形式化や自然言語への説明生成には依然として誤りが入りうる）。よって、「形式検証の及ぶ範囲を拡大する」、「問題文の形式化精度を向上させる」ことで生成AIの数学力は向上すると考えられており、研究が盛んに行われている。

    \end{body}

    \begin{tikzpicture}
        %========================
        % 共通
        %========================
        \pgfmathsetlengthmacro{\Span}{0.90*\linewidth}

        \pgfmathsetlengthmacro{\TopY}{0pt}
        \pgfmathsetlengthmacro{\BotY}{-3.5cm}

        \coordinate (TopL) at (-0.5*\Span,\TopY);
        \coordinate (TopR) at ( 0.5*\Span,\TopY);
        \coordinate (BotL) at (-0.5*\Span,\BotY);
        \coordinate (BotR) at ( 0.5*\Span,\BotY);

        %========================
        % 上段
        %========================
        \tikzset{
        box/.style={
        rectangle, draw, thick,
        text width=3.0cm,
        minimum height=1.0cm,
        align=center, font=\small,
        inner sep=2pt
        },
        myarrow/.style={->, >=Stealth, thick, shorten <=2pt, shorten >=2pt},
        danger/.style={box, fill=red!20},
        }

        \node[box] (nlProblem)   at ($(TopL)!0.00!(TopR)$) {自然言語\\問題};
        \node[box] (nlReasoning) at ($(TopL)!0.42!(TopR)$) {自然言語\\推論};
        \node[box] (nlAnswer)    at ($(TopL)!0.84!(TopR)$) {自然言語\\解答};

        \draw[myarrow] (nlProblem) -- (nlReasoning);
        \draw[myarrow] (nlReasoning) -- (nlAnswer);

        \node[above=3mm of nlProblem] {\bfseries 従来の自然言語推論};

        %========================
        % 下段
        %========================
        \tikzset{
        boxS/.style={
        rectangle, draw, thick,
        text width=1.85cm,
        minimum height=0.72cm,
        align=center,
        font=\scriptsize,
        inner sep=2pt
        },
        dangerS/.style={boxS, fill=red!20},
        safeS/.style={boxS, fill=green!20},
        myarrowS/.style={->, >=Stealth, thick, shorten <=2pt, shorten >=2pt}
        }

        \node[boxS]    (flProblemNl) at ($(BotL)!0.00!(BotR)$) {自然言語\\問題};
        \node[safeS]   (flProblem)   at ($(BotL)!0.21!(BotR)$) {形式言語\\問題};
        \node[safeS]   (flReasoning) at ($(BotL)!0.42!(BotR)$) {形式言語\\推論};
        \node[safeS]   (flAnswer)    at ($(BotL)!0.63!(BotR)$) {形式言語\\解答};
        \node[boxS] (finalAnswer) at ($(BotL)!0.84!(BotR)$) {自然言語\\解答};

        \draw[myarrowS] (flProblemNl) -- (flProblem);
        \draw[myarrowS] (flProblem)   -- (flReasoning);
        \draw[myarrowS] (flReasoning) -- (flAnswer);
        \draw[myarrowS] (flAnswer)    -- (finalAnswer);

        \begin{pgfonlayer}{background}
            \node[draw, thick, inner sep=2pt, rounded corners=2pt, fit=(flProblem)(flReasoning)(flAnswer), green, label={[font=\scriptsize]above:ハルシネーションを抑制}] (flBigBox) {};
        \end{pgfonlayer}

        \draw[myarrowS, dashed]
        (flReasoning.south) .. controls +(0,-0.8) and +(0,-0.8) ..
        node[pos=0.55, right, font=\tiny, xshift=-2mm, yshift=-2mm] {検証失敗}
        (flProblem.south);

        \node[above=3mm of flProblemNl] {\bfseries 形式検証統合推論};

    \end{tikzpicture}

    \begin{body}
        \indent 上記を踏まえると本稿の目的は、\emph{数学という体系（特にチェバの定理）を証明支援系を通してコンピュータに検証してもらうために適切な形式言語を選択しコンピュータに伝達すること}である。

        \indent この手段として私はLean4という言語を選択した。
    \end{body}

    \section{Lean4・mathlib4の概要}
    \begin{body}
        \indent Lean4とは、依存型理論の一種CIC（Constructive Inductive Calculus）を基礎理論とする、プログラミング言語であり証明支援系（ITP）でもある。

        \indent この言語の概要を知るにはまずLean4のファイル構成を把握するのが良いが、これは一般的な数学書と大きく変わらない。論理の流れに従い$def$、$lemma$や$theorem$という宣言に続いて用語の定義、命題の主張が並んでいる。数学書が一部の定理の証明を他の文献に譲るように、Lean4のファイルでは別のファイルを参照し、必要な定理をインポートすることでその定理を既知の定理として証明中で再利用できるようになる。再利用前提でファイルを作成する必要があるため、ファイル構成にもある程度の形式・ルールが生まれる。この要件を満たしながら定義・定理の数々をLean4に書き起こして行ったファイルを集めたものが、前述したmathlib4という巨大なライブラリである。証明支援系、ライブラリという人によっては馴染みのない用語を聞いて身構える読者もいるかもしれないが、Lean4を一種のプログラミング言語として捉えてしまえば、数学体系をLean4という言語で書き直し、mathlib4という共用のフォルダにまとめて管理しているだけである。世界中の数学者や技術者が共同でこのフォルダを管理・運営し、世界規模で開発が行われている。この言語は2021年にLean3を改良して生まれた新しい言語であるが、両者に互換性がないため、現在Lean3時代のライブラリ（mathlib3）からのLean4への翻訳が進められている。

        \indent そんなLean4であるが、数学書と最も大きく異なる点がその言語が型理論をベースにして構築されていることである。一般的な数学書は一階述語論理をベースに部分的に自然言語による説明を交えた形式で綴られるが、Lean4ではそのすべてが一つの言語を用いて記述される。数学的主張を世に出す立場として、慣れるまで少し時間がかかるのはデメリットであるが、一方でLean4が持つ大きなメリットが「証明を書く過程でその証明が正しいかどうかを判定してくれる点」である。背後で依存型理論をベースとした型チェッカーというツールが稼働し続け、記述している証明が誤っていれば直ちにエラーを返してくれる場合があり、非常に効率的に記述が行える。いわば、"添削機能付きの数学書"である。

        \indent Lean4における証明記述には大別して2つのモードが存在する。
        
        \begin{bodyitemize} [leftmargin=1.8em]
            \item[]
                \textbf{項モード}\par\noindent
                証明に相当する項をコーディングしながら自分で構成していく。
            \item[]
                \textbf{tacticモード}\par\noindent
                あらかじめ定義されたtacticと呼ばれるコマンドを利用して、証明すべき命題を変更するなどの操作を行う。
        \end{bodyitemize}

        \vskip\baselineskip

        \indent 詳細な説明は割愛するが、前者は証明全体の流れを明示しやすく，後者は手計算の証明手順に近い形で推論を整理できるため，両者は相補的に用いられる。
    \end{body}

    % =========================================================
    \chapter{構文論における定理の一般化}
    % =========================================================
    \begin{body}
        \indent この証明では一般化とはそもそも何か、定理の逆とは何かを構文論の観点から振り返る。
    \end{body}

    \section{定理の一般化とは}

    \begin{body}
        \indent 命題の集合$T$に対し、
        \[
        T \vdash \phi
        \]
        の形で表される命題$\phi$のことを（$T$から導出される）\emph{定理}と呼ぶ。特に$\phi$が$p\to q$の形のとき、命題$q\to p$を\emph{定理の逆}と呼ぶ。定理の逆は一般には真になるとは限らない。   

        \indent また、複数の命題をまとめて定理と呼ぶ場合も多い。$\phi_{1},\phi_{2},\cdots,\phi_{n}$が全て$T$から導出される定理である時、
        \[       
        T \vdash (\phi_{1}\land\phi_{2}\land\cdots\land\phi_{n})
        \]
        一方、自然演繹を推論規則として採用する上では演繹定理と呼ばれる以下の性質が成り立つ。
        \[
        T\cup\{\psi\}\vdash \phi \;\Longleftrightarrow\; T\vdash (\psi\to\phi).
        \]
        $T= \{\varphi_{1},\varphi_{2},\cdots,\varphi_{n}\}\cup\{\psi_{1},\psi_{2},\ldots,\psi_{n}\}$の時、
        \[
        \begin{alignedat}{2}
        &\phantom{\iff}\; T \vdash \phi && \\
        &\iff \{\varphi_{1},\varphi_{2},\cdots,\varphi_{n}\}\vdash\bigl(\psi_{1}\to(\psi_{2}\to(\cdots\to(\psi_{n}\to\phi)))\bigr) &&\text{\footnotesize（演繹定理を繰り返し適用）}\\
        &\iff \{\varphi_{1},\varphi_{2},\cdots,\varphi_{n}\}\vdash\bigl(\neg\psi_{1}\lor(\neg\psi_{2}\lor(\cdots\lor(\neg\psi_{n}\lor\phi)))\bigr) &&\text{\footnotesize（含意と論理和の関係）}\\
        &\iff \{\varphi_{1},\varphi_{2},\cdots,\varphi_{n}\}\vdash\bigl(\neg\psi_{1}\lor\neg\psi_{2}\lor\cdots\lor\neg\psi_{n}\lor\phi\bigr) &&\text{\footnotesize（$\lor$の結合法則）}\\
        &\iff \{\varphi_{1},\varphi_{2},\cdots,\varphi_{n}\}\vdash\bigl(\neg(\psi_{1}\land\psi_{2}\land\cdots\land\psi_{n})\lor\phi\bigr) &&\text{\footnotesize（ド・モルガンの法則）}\\
        &\iff \{\varphi_{1},\varphi_{2},\cdots,\varphi_{n}\}\vdash\bigl((\psi_{1}\land\psi_{2}\land\cdots\land\psi_{n})\to\phi\bigr) &&\text{\footnotesize（含意と論理和の関係）}\\
        \end{alignedat}
        \]
        よって、$\phi_{1},\phi_{2},\cdots,\phi_{n}$が$T$から導出される定理である時、以下のように書き下せる。
        \[
        \{\varphi_{1},\varphi_{2},\cdots,\varphi_{n}\}\vdash{\psi_{1}\land\psi_{2}\land\cdots\land\psi_{n}}\to{\phi_{1}\land\phi_{2}\land\cdots\land\phi_{n}}
        \]
        
        \vskip\baselineskip

        $P:=\{\varphi_{1},\varphi_{2},\cdots,\varphi_{n}\},A:=\{\psi_{1},\psi_{2},\cdots,\psi_{n}\},C:=\{\phi_{1},\phi_{2},\cdots,\phi_{n}\}$と置き、便宜上以下のように略記する
        \[
        P\vdash A\to C
        \]
        これは自然演繹の公式の記法ではなく、説明上の都合でこのような書き方になっていることに留意してほしい。
        この$P$を\emph{前提/背景理論}、$A$を\emph{仮定}、$C$を\emph{結論}と呼ぶ。
        要するに、定理が与えられるとそれを適当な粒度の命題の組を用意し、\emph{前提・仮定・結論の導出関係として書き換えることができる}ということである。

        \indent 数学基礎論以外の分野の一般的な数学書では前提$P$として$ZF/ZFC$が慣習的に用いられるため、これらは暗黙に受け入れたものとして敢えては明記されず、仮定$A$から結論$C$を導く、というステップのみが記述されることが多い。

        \indent また、この形式で定理を表現することにより、暗黙に認めていた前提を明文化できるため、Lean4での形式化に繋げやすいというメリットがある。

        \indent この上で、本稿で言う「定理の一般化」とは、既存の定理がある意味で\emph{特別な場合として含まれる}ような、より広い主張をする定理を与えることである。典型的には次の三方向がある。
        \begin{bodyitemize}[leftmargin=3.5em]
            \item \textbf{前提の弱化}：$P$の元の命題の主張を弱める。$\varphi\in P$を、$\varphi\to \varphi'$を満たす$\varphi'$で置き換える。 
            \item \textbf{仮定の弱化}：$A$の元の命題の主張を弱める。$\psi\in A$を、$\psi\to \psi'$を満たす$\psi'$で置き換える。
            \item \textbf{結論の強化}：$C$の元の命題の主張を強める。$\phi\in C$を、$\phi'\to \phi$を満たす$\phi'$で置き換える。
            \item \textbf{複数の組合せ}：上記を複数個組み合わせる。
        \end{bodyitemize}
        
        \indent (3)は「定理の強化」と呼べれることが多いが、本稿では一般化に含める。一般化と聞くと仮定と結論に着目しがちであるが、定理が依存する前提を緩め（例：体から環へ、可換環から一般の環へ、有限次元から任意次元へ）、対象の主張をより広い土台で成立させることも一般化の重要な形である。

    \end{body}

    % =========================================================
    \chapter{チェバの定理とアフィン空間}
    % =========================================================
    \begin{body}
        \indent チェバの定理とは高校数学でも頻繁に登場する、初等幾何における重要な定理の一つである。

        \indent 本章ではこのチェバの定理の詳細と、チェバの定理が展開されるアフィン空間の性質について議論する。

    \end{body}

    \section{チェバの定理の有名な形}

    \begin{body}
        \indent 高校の教育課程において、チェバの定理は以下の形で紹介されることが多い。

        \indent ただし、定理中で登場する用語は初等幾何的な高校の学習課程内の意味である。
    \end{body}

    \begin{theorem}[古典的チェバの定理]\label{thm:classical-ceva}
        三角形$\triangle ABC$の各辺$BC,CA,AB$またはその延長線上にそれぞれ頂点とは異なる点$D,E,F$をとる。
        このとき$3$直線$AD$、$BE$、$CF$が$1$点で交わるならば、以下の式が成立する
        \[
          \dfrac{BD}{DC}\,\dfrac{CE}{EA}\,\dfrac{AF}{FB} = 1
        \]
    \end{theorem}

    \begin{theorem}[古典的チェバの定理の逆]
        三角形$\triangle ABC$の各辺$BC,CA,AB$またはその延長線上にそれぞれ頂点とは異なる点$D,E,F$をとる。
        このとき
        \[
          \dfrac{BD}{DC}\,\dfrac{CE}{EA}\,\dfrac{AF}{FB} = 1
        \]
        が成立するならば、$3$直線$AD,BE,CF$は$1$点で交わる、または互いに平行である。
    \end{theorem}

    \begin{body}
        \indent ここで登場する点$D,E,F$を\emph{チェビアンの足}、$3$直線$AD,BE,CF$を\emph{チェビアン}と呼ぶ。

        \indent つまり古典的チェバの定理（またはその逆）とは\emph{チェビアンが共点するという幾何的性質が成立する条件を代数的に記述する定理}である。

        \indent 上記のチェバの定理の主張には、一見すると線分の長さ（BD、DC等）が現れるため、ユークリッド幾何に依存する定理のように見える。

        \indent しかし、実際にはこの定理の主張は線分の長さの比（$\dfrac{BD}{DC}$ 等）のみで記述されている。実はこれらの概念は後述するようにユークリッド幾何を用いることなくアフィン幾何という体系内で定義することができる。

        \indent またチェバの定理が必要十分性を主張する定理だと認識していた読者も多いかもしれない。$D,E,F$を辺上に限定する版や、$3$直線が$1$点で交わる必要十分条件を提供する版など、教科書により紹介の仕方に微妙な違いがあり、これらは教育的配慮によるものである。詳細は\hyperref[appendix:classical-ceva-version]{付録}で解説するが、アフィン幾何の範疇では古典的チェバの定理の主張したい内容が「共点$⇔$等式条件」というシンプルな形の必要十分条件では表せないという事にだけ留意してほしい。

        \indent 以降、アフィン幾何とはどのような体系なのか紹介する。
    \end{body}

    
    \section{概念の定義}\label{sec:affine-prelim}
    \begin{body}
        \indent 以降の議論で必要な概念だけを、\emph{証明で頻繁に使う順}に導入する。
        本稿の主眼はチェバ型定理の形式化であるため、凸性や距離など計量的な概念は一切用いない。
        （したがって順序環やノルム空間の仮定も不要である。）

        \indent 以下、特に断らない限り $k$ は環、$V$ は $k$-加群、$P$ は $V$ 上のアフィン空間、$\iota$ は添字集合とする。
        また、有限和は $\sum_{i\in s}$ の形で書くが、Lean/mathlib4 では対応する概念は \texttt{Finset} 上の和である。
    \end{body}

    \begin{definition}[アフィン空間（mathlib: \texttt{AffineSpace V P}）]
        $P$ が $V$ 上の（$k$ に関する）\emph{アフィン空間}であるとは、
        \[
          \dotplus : V\times P \to P,\qquad \dotminus : P\times P \to V
        \]
        が与えられ、任意の $p,q\in P$ と $v\in V$ について
        \[
          \mathbf 0 \dotplus p = p,\qquad (p\dotminus q)\dotplus q = p,\qquad (v\dotplus p)\dotminus p = v
        \]
        が成り立つことをいう。$P$ の元を\emph{点}と呼ぶ。
        mathlib4 ではこの構造は \texttt{AddTorsor V P} として定義され、\texttt{open scoped Affine} を行うと \texttt{AffineSpace V P} と表示される。
    \end{definition}

    \begin{definition}[アフィン結合（mathlib: \texttt{Finset.affineCombination}）]\label{def:affcomb}
        $s\subseteq \iota$ を有限集合とし、重み $\lambda:s\to k$ が
        \[
          \sum_{i\in s}\lambda(i)=1
        \]
        を満たすとき、任意に $j\in s$ を固定して
        \[
          \AffComb(p|_s,\lambda)
          \;:=\;
          \left(\sum_{i\in s}\lambda(i)\,\bigl(p(i)\dotminus p(j)\bigr)\right)\dotplus p(j)
        \]
        と定める。これを点族 $p|_s$ の\emph{アフィン結合}と呼び、$\lambda$ を\emph{重み}と呼ぶ。
    \end{definition}

    \begin{remark}[基点に依らないことと略記]
        条件 $\sum_{i\in s}\lambda(i)=1$ のもとで、上の定義は $j\in s$ の選び方に依らない。
        以降、非公式な略記として
        \[
          \AffComb(p|_s,\lambda)=\sum_{i\in s}\lambda(i)\,p(i)
        \]
        と書くことがある（右辺は「点の線形結合」という本来未定義の演算を含む略記である）。
        Lean では \texttt{Finset} を用い、$s$ を有限集合そのものとして扱って \texttt{s.affineCombination k p w} と書く。
    \end{remark}

    \begin{definition}[アフィン包（mathlib: \texttt{affineSpan}）]
        点族 $p|_s$ の\emph{アフィン包}を
        \[
          \AffHull(p,s):=
          \left\{\,
            \AffComb(p|_s,\lambda)\ \middle|\ \lambda:s\to k,\ \sum_{i\in s}\lambda(i)=1
          \,\right\}\subseteq P
        \]
        で定める。mathlib4 の \texttt{affineSpan k (p '' s)} の台集合（\texttt{carrier}）に対応する。
    \end{definition}

    \begin{definition}[直線と平行（mathlib: \texttt{$line[k, p_1, p_2$]}, \texttt{$AffineSubspace.Parallel$}）]
        2点 $A,B\in P$ に対し、2点集合 $\{A,B\}$ のアフィンスパンを\emph{直線}と呼び、
        \[
          \mathrm{line}[k,A,B]:=\mathrm{affineSpan}_k(\{A,B\})
        \]
        と書く（Lean では \texttt{line[k, A, B]}）。
        また 2 つのアフィン部分空間（特に直線）が平行であることを $\parallel$ で表し、Lean では \texttt{∥}（\texttt{AffineSubspace.Parallel}）で表す。
    \end{definition}


    \begin{definition}[アフィン独立（mathlib: \texttt{AffineIndependent}）]
        点族 $p:\iota\to P$ が\emph{アフィン独立}であるとは、
        任意の有限集合 $s\subseteq \iota$ と重み $\lambda,\mu:s\to k$ について
        \[
          \sum_{i\in s}\lambda(i)=1,\ \sum_{i\in s}\mu(i)=1,\ 
          \AffComb(p|_s,\lambda)=\AffComb(p|_s,\mu)
          \ \Longrightarrow\ 
          \lambda=\mu
        \]
        が成り立つことをいう。
        直観的には「重心座標が一意に定まる」という性質である。
    \end{definition}

    \begin{definition}[単体と面（mathlib: \texttt{Affine.Simplex}, \texttt{face}, \texttt{faceOpposite}）]
        $n\in\mathbb N$ に対し、$n+1$ 個の点からなる点族 $p:\mathrm{Fin}(n+1)\to P$ がアフィン独立であるとき、
        これを\emph{$n$-単体}（simplex）と呼ぶ。
        特に $n=2$ のとき $2$-単体は三角形であり、Lean では \texttt{Triangle} と略記されている。

        \indent また、頂点集合の部分集合（\texttt{Finset}）を指定することで得られる部分単体を\emph{面}（face）と呼び、
        頂点 $i$ を除いた面を\emph{対向面}（\texttt{faceOpposite}）と呼ぶ。
        本稿ではチェバ型定理の一般化に必要な範囲でのみこれらを用いる。
    \end{definition}

    \begin{definition}[共点]
        直線 $L_1,\dots,L_m$ が\emph{共点する}とは、
        ある点 $x\in P$ が存在して $x\in\bigcap_{j=1}^m L_j$ が成り立つことをいう。
        この $x$ を\emph{共点}と呼ぶ。
    \end{definition}

    \section{チェバの定理のアフィン性}\label{sec:ceva-affine}
    \begin{body}
        \indent 古典的チェバの定理は一見すると線分の長さが現れるためユークリッド幾何の定理に見える。
        しかし本質は「\emph{同一直線上の 2 点を端点とするアフィン結合}」と「\emph{直線の共点性}」だけで記述できる。

        \indent 実際、辺 $BC$ 上（または延長上）の点 $D$ は、ある係数 $t\in k$ を用いて
        \[
          D = (1-t)B + tC
        \]
        （すなわち $B,C$ のアフィン結合）として表せる。高校幾何で現れる比 $\frac{BD}{DC}$ は、
        この $t$ に対し（有向線分の意味で）
        \[
          \frac{BD}{DC} = \frac{t}{1-t}
        \]
        と対応する。よって古典的チェバの等式条件は「係数 $t,u,v$ による条件」として書ける。

        \indent Lean/mathlib4 では、この対応をさらに機械的に扱えるよう、
        2 点間のアフィン結合を \texttt{AffineMap.lineMap}、
        直線を \texttt{line[k, -, -]}（2 点のアフィンスパン）で表現する。
        その上で、チェバの定理（順方向）は \texttt{Mathlib.LinearAlgebra.AffineSpace.Ceva} において、
        \emph{アフィン独立性とアフィン結合のみに依存する定理}として証明されている。
    \end{body}

    \begin{theorem}[mathlib4 における三角形のチェバ（順方向）{\small（\texttt{Affine.Triangle.prod\_eq\_prod\_one\_sub\_of\_mem\_line\_point\_lineMap}）}]
        $k$ を可換環（零因子を持たない）とし、三角形 $t$ と係数 $r:\mathrm{Fin}\,3\to k$ を与える。
        ある点 $p'\in P$ が各頂点 $t.\mathrm{points}\,i$ と対向辺上の点
        $\mathrm{lineMap}(t.\mathrm{points}(i+1))(t.\mathrm{points}(i+2))(r(i))$
        を結ぶ直線上にあるとき、
        \[
          \prod_{i:\mathrm{Fin}\,3} r(i)\;=\;\prod_{i:\mathrm{Fin}\,3}(1-r(i))
        \]
        が成り立つ。
    \end{theorem}

    \begin{proof}[証明の概略（なぜ「アフィン幾何の結果」なのか）]
        mathlib4 の証明は、まず一般添字版のチェバ型定理
        \[
          \texttt{AffineIndependent.exists\_affineCombination\_eq\_smul\_eq}
        \]
        を示し、ついで $ι=\mathrm{Fin}\,3$ の特殊化として上の三角形版を導く、という構造になっている。
        一般添字版は \emph{加群・アフィン空間・アフィン結合・アフィン独立}しか仮定しておらず、
        距離・角度・内積などの計量構造は一切用いていない。
        したがって三角形版も（環上の）アフィン幾何の結果である。
        詳細は \texttt{Mathlib/LinearAlgebra/AffineSpace/Ceva.lean} の実装に従う。
    \end{proof}

    \begin{body}
        \indent 以降の章では、この「一般添字版チェバ型定理」を出発点として、
        その\emph{逆方向}（与えられた係数整合条件から共点性を導く）を、
        mathlib4 に適した形で追加することを目標とする。
    \end{body}


    % =========================================================
    \chapter{チェバの定理の一般化と選定}
    % =========================================================
    \begin{body}
        \indent 本章では上記の議論を踏まえ、一般化チェバの定理の主張として求められる何が相応しいかを明確にする。
    \end{body}

    \section{一般化の方針}

    \begin{body}
        \indent 古典的チェバに対して上記の定式化を行うと定理が様々な観点から一般化できることが見える。

        \indent 例えば以下の観点が考えられる。

        \begin{bodyitemize}
            \item 周辺空間のみ一般化する
            \item $2$-単体を$n$-単体に一般化する
            \item $1$-チェビアンを$2$-チェビアンで考えてみる
            \item チェビアンの数は頂点数より多くとも構わない
        \end{bodyitemize}

        \vskip\baselineskip

        \indent 実際に一般化したものとして以下のようなものが挙げられる。

        \begin{bodydescription}
            \item[\hspace*{1.5em}・候補A：周囲空間のみ一般化]
            $n$次元アフィン空間に三角形を埋め込み、その各頂点から対向する$1$-面上の点へ下ろす$1$-チェビアンを考える．
            これらの直線が（$n$次元アフィン空間内で）共点（または平行）であることと、古典的チェバと同型の比の積等式が成り立つことの必要十分性を主張する定理．
            （本質は古典的チェバの定理と同じで、周囲空間の次元は主張に影響しない．）
            
            \item[\hspace*{1.5em}・候補B：$n$-単体版, 1-チェビアン]
            $n$-単体の各頂点から対向する$(n-1)$-面上の点へ下ろす$1$-チェビアン（直線）を考える．
            これらが共点または平行であることと，同次質量分布型の条件（頂点への質量割当から各対向面上の点が重心として与えられる等）が成り立つことの必要十分性を主張する定理．
            
            \item[\hspace*{1.5em}・候補C：$n$-単体版, 2-チェビアン]
            $n$-単体の各頂点から対向する$(n-1)$-面に対し，$2$次元的データ（面積比・外積・行列式など）で自然に定まる$2$-チェビアンを考える．
            これらが共点または平行であることと，面積比（あるいは行列式）により与えられる条件が成り立つことの必要十分性を主張する定理．
            
            \item[\hspace*{1.5em}・候補D：$n$単体版, multipede]\cite{multipede}
            $n$-単体の各頂点から対向する$k$-面（$1 \leq k \leq n-1$）への$1$-チェビアンを考える（$k$を固定しても，全ての$k$を同時に扱ってもよい）．
            これらが所定の意味で共点（または平行）であることと，対応する係数整合（バリセントリック係数の整合性，または質量分布条件の一般化）が成り立つことの必要十分性を主張する定理．
            一つの頂点から複数のチェビアンが伸びることから$multipede$版チェバと呼ばれる。
                
        \end{bodydescription}

        \indent このような、さまざまな版で一般化されたチェバの定理の総称を\emph{チェバ型定理}と呼ぶ。
    \end{body}

    \section{先行プロジェクトの紹介}
    \begin{body}
        \indent 過去のいくつかのチェバ型定理形式化に関する先行研究を紹介する。

        \begin{bodydescription}
            \item[\hspace*{1.5em}・Lean3時代の形式化]
            期間: $2021$年$12$月 〜 $2023$年$7$月\\
            GitHub: \url{https://github.com/leanprover-community/mathlib3/pull/10632} \cite{mathlib3_ceva}

            Mantas Bakšys氏による試み。
            $2$-単体のチェバの定理のみを形式化、ユークリッド幾何の概念である距離を全面に出した証明を行ったが、mathlib4に追加するには一般性が足りないとして採用されなかった。

            \item[\hspace*{1.5em}・Aristotle.AI（Harmonic社）を用いた形式化]
            期間: $2025$年$12$月 〜 $2026$年$1$月\\
            GitHub: \url{https://github.com/leanprover-community/mathlib4/pull/33388} \cite{mathlib4_ceva_alexeev}

            Boris Alexeev氏による試み。
            Harmonic社が提供するLeanプログラム特化のAIサービスAristotle.AIを用いて形式化した。
            全てのコードがLeanの機械検証を通過し、正確性が保証された。AIの性能向上を示したと言える。
            しかし、$2$-単体のチェバの定理のみを形式化しており、またコミュニティ内の過去の議論を踏まえられていなかったため、mathlib4に追加するには一般性が足りないとして採用されなかった。

            \item[\hspace*{1.5em}・$n$-単体で汎用一般化された初の形式化]
            期間: $2026$年$1$月 〜 $2026$年$1$月\\
            GitHub: \url{https://github.com/leanprover-community/mathlib4/pull/33409} \cite{mathlib4_ceva_myers}

            Joseph Myers氏による試み。
            過去の不十分な事例を踏まえて、$n$-単体に一般化されたチェバの定理順方向を形式化した。
            おそらく現時点で最も一般性が高い形式化である。
            後の章で詳細に見ていく。
            2026/1/14 PR受理。
        \end{bodydescription}

    \end{body}

    \section{本稿で扱う一般化の選定}
    \begin{body}
        \indent 本稿では既存のmathlib4に用意された一般化チェバの定理の順方向の方針を採用する。

        \indent 上述したJoseph Myers氏によって形式化されたものであり、「各頂点と、その対向面または延長上の両者を結ぶ$1$-チェビアンが共点する必要条件」を与える。
    \end{body}

    \subsection{先行研究の主張}

    \begin{body}
        \indent mathlib4内のCeva.leanは主に以下の定理が証明されている。
        
        \begin{bodyenumerate}
            \item (2)を示すための補題
            \item 任意添字点族チェバの定理
            \item 有限添字点族チェバの定理
            \item $2$-単体についての古典的チェバの定理（積等式形）
            \item $2$-単体についての古典的チェバの定理（商等式形）
        \end{bodyenumerate}
        \indent 本稿のテーマに近いのは(3)である。

        \indent (1)で一般化チェバの定理の主張のほとんどを補題として提示、(2)で任意添字点族に対する一般化チェバの定理を示し、その特殊な版として(3)を導いている。
            
        \indent (4),(5)ではさらに(3)の特殊形として3つの添字に対する古典的チェバの定理を積等式と商等式で版を分けて紹介している。

        \indent (3)の主張を敢えて要約すると以下のようになる。
    \end{body}

    \begin{theorem}[チェバ型定理順方向]
        $n$-単体の頂点を$p(0)$,…,$p(n)$とする。
        いくつかの頂点$p(i)$に対し、$p(i)$を含まない任意の面上に点$q(i)$を固定し、直線$p(i)q(i)$を考える。
        こうして得られる複数の直線がある点$p'$で交わる時、$p'$は$p(F_i)\cup \{p(i)\}$が張るアフィン部分空間中に存在して、各iについて$q(i)$の$F_i$が定める点族に関する重心座標、および$p'$の$\iota$が定める点族に関する重心座標は第$i$成分を除いて比例する。
    \end{theorem}

    \begin{body}
        \indent もはや古典的チェバの定理の原形を留めていない。

        \indent 点$q(i)$、直線$p(i)q(i)$それぞれは古典的チェバの定理においてチェビアンの足だったもの、チェビアンだったものである。要するに、「共点が存在すれば、共点とチェビアンの足の重心座標成分の大部分が比例する」ことを主張している。

        \indent 一見古典的チェバとの関連は見えづらい。しかし、以下の点で一般化になっている。

        \indent 一つ目に、\emph{$n$-単体についての主張になっている}。前提の弱化である。$n=2$とすれば三角形についての定理になる。

        \indent 二つ目に、\emph{チェビアンは全ての頂点から下ろす必要がない}。前提の弱化である。直線$p(i)q(i)$がチェビアンに相当する。古典的チェバの定理では、全ての頂点からチェビアンを下ろしている。

        \indent 三つ目に、\emph{チェビアンの足は対向面上でなくて良い}。前提の弱化である。点$q(i)$がチェビアンの足に相当する。古典的チェバの定理では、チェビアンの足は対向面（対辺）上にあった。しかし、この定理中ではチェビアンの足はいずれかの面上にあれば良い。

        \indent 四つ目に、\emph{商の等式条件が係数の比例条件になっている}。結論の強化である。これによりスカラーを体の元に限る必要がなくなる。

        \indent この主張を導出関係を明確にして書き下す。

    \end{body}

    \begin{theorem}[チェバ型定理順方向の言い換え]
        \noindent\textbf{\\前提：}\par
        \begin{theoremenumerate}[label=(P\arabic*), leftmargin=4em, topsep=2pt, partopsep=0pt, itemsep=0pt, parsep=0pt]
            \item $k$ は 環である。
            \item $V$ は $k$-加群である。
            \item $P$ は $V$ 上のアフィン空間である。
            \item $\iota$ は有限の添字集合である。
            \item 点族 $p : \iota \to P$はアフィン独立。
            \item $s \subseteq \iota$ は空でない。
            \item 各 $i\in s$ について，有限集合 $F_i\subseteq \iota$ と重み $w_i:F_i\to k$ が与えられ，
            \[
            i\notin F_i \land \sum_{j\in F_i} w_i(j)=1
            \]
            が成り立つ。
        \end{theoremenumerate}
        \noindent\textbf{\\仮定：}\par
        \begin{theoremenumerate}[label=(A\arabic*), leftmargin=4em, topsep=2pt, partopsep=0pt, itemsep=0pt, parsep=0pt]
            \item ある点 $p'\in P$ が存在し、各 $i \in s$ について、$p'$ は2点 $p(i)$ と $\sum_{j \in F_i} w_i(j)p(j)$ を結ぶ直線上に存在する。
        \end{theoremenumerate}

        \noindent\textbf{\\結論：}\par
        \begin{theoremenumerate}[label=(C\arabic*), leftmargin=4em, topsep=2pt, partopsep=0pt, itemsep=0pt, parsep=0pt]
            \item ある重み $w' : \iota \to k$が存在して
            \[
            \sum_{j \in \iota} w'(j) = 1
            \]
            を満たし、
            \[
            p' = \sum_{j \in \iota} w'(j)p(j)
            \]
            と書ける。
            \item 各 $i \in s$ ごとにスカラー $r_i \in k$ が存在し、次が成り立つ： 
            \[
            \forall i \in s,\ \exists r_i \in k,\ \forall j \in F_i,\ \bigl(w'(j) = r_i\, w_i(j)\bigr).
            \]
        \end{theoremenumerate}
    \end{theorem}

    \begin{proof}  
        (A1) より、任意の $i\in s$ に対して $p'$ は2点 $p(i)$ と $\sum_{j \in F_i} w_i(j)p(j)$ を結ぶ直線上に存在する。
        よって、
        \[
        \exists \alpha_i\in k, p'=\alpha_i\,p(i)+(1-\alpha_i)\sum_{j\in F_i} w_i(j)\,p(j)
        \]
        ここで各 $i\in s$ について、関数 $\bar w_i:\iota\to k$ を次で定める：
        \[
        \bar w_i(j):=
        \begin{cases}
        \alpha_i & (j=i),\\
        (1-\alpha_i)\,w_i(j) & (j\in F_i),\\
        0 & (j\in\iota\setminus(F_i\cup\{i\})).
        \end{cases}
        \]
        （$w_i: F_i\to k$ であり、$\iota$全域で定義されていないため、$j\notin F_i$ の場合を上の定義により$\bar w_i(j)=0$ として埋めた。）
  
        $i\notin F_i$ と $\sum_{j\in F_i}w_i(j)=1$ より
        \[
        \begin{aligned}
        \sum_{j\in\iota} \bar w_i(j)
            &= \bar w_i(i)+\sum_{j\in F_i} \bar w_i(j)+\sum_{j\in\iota\setminus(F_i\cup\{i\})} \bar w_i(j)\\
            &= \alpha_i + (1-\alpha_i)\sum_{j\in F_i} w_i(j) + 0\\
            &= \alpha_i + (1-\alpha_i)\\
            &= 1
        \end{aligned}
        \]
        が成り立つ。また上の定義から
        \[
        p'=\sum_{j\in\iota} \bar w_i(j)\,p(j)
        \]
  
        以上より、各 $i\in s$ に対して
        \[
        \sum_{j\in\iota} \bar w_i(j)=1,\qquad
        p'=\sum_{j\in\iota} \bar w_i(j)\,p(j)
        \]
        が得られた。
  
        また(P5)より、$p:\iota\to P$はアフィン独立であるから$\AffHull(p)$中の点を点族$p$のアフィン結合で一意に表せる。
        (A1)より、$p'\in P$である。
        従って任意の $i,i'\in s$ について
        \[
        \bar w_i = \bar w_{i'}
        \]
  
        よって、任意に $i_0\in s$ を一つ固定して
        \[
        w' := \bar w_{i_0}:\iota\to k
        \]
        とおけば、全ての $i\in s$ に対して $w'=\bar w_i$。
  
        したがって
        \[
        \sum_{j\in\iota} w'(j)=1,\qquad
        p'=\sum_{j\in\iota} w'(j)\,p(j)
        \]
        (C1) が示された。
  
        また、$r_i:=1-\alpha_i$ とおけば、定義より$w'=\bar w_i$ なので、任意の $j\in F_i$ について
        \[
        \begin{aligned}
            \forall j\in F_i, 
            w'(j) &=\bar w_i(j)
                &=(1-\alpha_i)\,w_i(j)
                &=r_i\,w_i(j)
        \end{aligned}
        \]
        (C2)が示された。
    \end{proof}

    \begin{body}
        \indent これが古典的チェバの定理の一般化になっていることは、前提・仮定・結論に含まれる命題の強弱を比較すれば直ちにわかる。

        \indent mathlib4内のCeva.Leanでは上記の主張をベースにLean4向けに整形した証明がなされている。

        \indent この定理を参考に、本稿では古典的チェバの定理の逆を一般化して示し形式化する。
    \end{body}

    % =========================================================
    
    % =========================================================
    \chapter{Lean4での実装}\label{ch:implementation}
    % =========================================================

    \section{目標：一般化チェバ型定理の「逆方向」}\label{sec:goal}
    \begin{body}
        \indent 本研究が直接追加したいのは、\texttt{Mathlib.LinearAlgebra.AffineSpace.Ceva} にある一般添字版チェバ型定理
        \[
          \texttt{AffineIndependent.exists\_affineCombination\_eq\_smul\_eq}
        \]
        の「逆向き」である。

        \indent 注意すべき点は二つある。
        一つ目に、古典的チェバの「逆」（等式条件 $\Rightarrow$ 共点 \emph{または} 平行）は、
        射影幾何（無限遠点）を暗に含んだ主張であり、mathlib4 の既存ファイルが採用している
        \emph{純アフィン幾何}の枠内ではそのまま「$\Leftrightarrow$」にできない。
        二つ目に、Myers の一般化定理自体が「共点性 $\Rightarrow$ 係数整合」という\emph{必要条件}の形で書かれているため、
        ここでいう「逆方向」とは\emph{その結論（係数整合）から元の共点性仮定を回収する}という意味である。

        \indent したがって本章では、まず \texttt{exists\_affineCombination\_eq\_smul\_eq} の結論を仮定に持つ補題を追加し、
        「結論 $\Rightarrow$ 仮定」を Lean で証明する。これにより、少なくとも
        \[
          \text{（共点性）} \ \Longleftrightarrow\ \text{（重みの整合性）}
        \]
        という\emph{同値形}を、必要に応じてユーザ側で組み立てられるようになる。
    \end{body}

    \section{目標定理（自然言語）}\label{sec:goal-statement}
    \begin{theorem}[一般添字チェバ型定理（逆方向：回収補題）]\label{thm:ceva-converse-target}
        \noindent\textbf{前提：}\par
        \begin{theoremenumerate}[label=(P\arabic*), leftmargin=4em, topsep=2pt, partopsep=0pt, itemsep=0pt, parsep=0pt]
            \item $k$ は環、$V$ は $k$-加群、$P$ は $V$ 上のアフィン空間である。
            \item 点族 $p:\iota\to P$ はアフィン独立である。
            \item $s\subseteq \iota$ は空でない。
            \item 各 $i\in s$ について有限集合（\texttt{Finset}） $fs(i)\subseteq \iota$ と重み $w(i):\iota\to k$ が与えられ、
            \[
              \sum_{j\in fs(i)} w(i)(j)=1
            \]
            が成り立つ。
        \end{theoremenumerate}

        \noindent\textbf{仮定：}\par
        \begin{theoremenumerate}[label=(A\arabic*), leftmargin=4em, topsep=2pt, partopsep=0pt, itemsep=0pt, parsep=0pt]
            \item ある重み $w':\iota\to k$ と有限集合 $fs'\subseteq \iota$ が存在して
            \[
              \sum_{j\in fs'} w'(j)=1,\qquad fs'.\AffComb(p,w')=p'
            \]
            を満たす（ここで右辺は \texttt{fs'.affineCombination k p w'}）。
            \item 各 $i\in s$ についてある $r_i\in k$ が存在して、任意の $j\in \iota$ について
            \[
              r_i\cdot \mathbf{1}_{fs(i)\setminus\{i\}}(j)\,w(i)(j)
              \;=\;
              \mathbf{1}_{fs'\setminus\{i\}}(j)\,w'(j)
            \]
            が成り立つ（$\mathbf 1$ は指示関数で、Lean では \texttt{Set.indicator}）。
        \end{theoremenumerate}

        \noindent\textbf{結論：}\par
        \begin{theoremenumerate}[label=(C\arabic*), leftmargin=4em, topsep=2pt, partopsep=0pt, itemsep=0pt, parsep=0pt]
            \item 任意の $i\in s$ について
            \[
              p' \in \mathrm{line}\!\left[k,\ p(i),\ (fs(i)).\AffComb(p,w(i))\right]
            \]
            が成り立つ。すなわち $p'$ は各「頂点 $p(i)$ と面上の点 $(fs(i)).\AffComb(p,w(i))$ を結ぶ直線」上にある。
        \end{theoremenumerate}
    \end{theorem}

    \section{証明方針（紙面）}\label{sec:proof-idea}
    \begin{body}
        \indent 証明の鍵は次の一点である：\emph{アフィン独立性は重心座標の一意性を保証する}。
        よって、同一点 $p'$ を二通りのアフィン結合で表したとき、重みは（有限集合の外側を 0 とみなせば）一致する。

        \indent 具体的には、各 $i\in s$ について「$p(i)$ と $(fs(i)).\AffComb(p,w(i))$ を結ぶ直線上の点」は、
        ある係数 $t_i\in k$ により
        \[
          \mathrm{lineMap}\bigl(p(i)\bigr)\bigl((fs(i)).\AffComb(p,w(i))\bigr)(t_i)
        \]
        と書ける。これを全頂点 $p$ のアフィン結合に展開すると、
        $j\neq i$ の成分は $t_i$ 倍で一様にスケールされる。
        Myers の順方向定理の結論がまさにこの「スケール不変量」を指示関数で抜き出した等式である。

        \indent 逆方向では、(A2) の指示関数等式が与えられているので、
        それを用いて「$p'$ の重み $w'$ が、ある $t_i$ による \texttt{lineMap} 展開と一致する」ことを示す。
        最後に \texttt{mem\_affineSpan\_pair\_iff\_exists\_lineMap\_eq} 型の補題により
        $p'$ の直線所属が従う。
    \end{body}

    \section{mathlib4 の既存実装の読み取り}\label{sec:read-mathlib}
    \begin{body}
        \indent \texttt{Mathlib.LinearAlgebra.AffineSpace.Ceva}（\texttt{Ceva.lean}）は、大きく次の三層構造になっている。
        \begin{bodyenumerate}
            \item \textbf{一般添字版（\texttt{AffineIndependent} 名前空間）}\\
            \texttt{exists\_affineCombination\_eq\_smul\_eq}：
            共点仮定から「グローバルな重み $w'$ の存在」と「指示関数で切り出した成分の比例」を得る。
            \item \textbf{有限添字版（\texttt{Fintype} 仮定）}\\
            \texttt{exists\_affineCombination\_eq\_smul\_eq\_of\_fintype}：
            $fs(i)=\texttt{Finset.univ}$ に固定して書きやすくした版。
            \item \textbf{三角形への特殊化（\texttt{Affine.Triangle} 名前空間）}\\
            係数 $r:\mathrm{Fin}\,3\to k$ と \texttt{lineMap} を用いて、積等式版・除法版を導く。
        \end{bodyenumerate}

        \indent 特に一般添字版の証明では、直線所属を
        \[
          p' \in \texttt{line[k, p i, ...]}
        \]
        から
        \[
          \exists r,\ \texttt{AffineMap.lineMap ... r = p'}
        \]
        に言い換え、係数 $r$ を「重み関数の \texttt{lineMap}」に持ち上げて
        指示関数等式を導いている。逆方向はこの流れをほぼ逆にたどるのが自然である。
    \end{body}

    \section{実装設計：追加すべき補題の最小集合}\label{sec:design}
    \begin{body}
        \indent 本稿の目的に照らすと、不要な概念や述語を増やさず、既存の補題で繰り返し出てくる形を
        「一回だけ名前を付ける」方針が望ましい。今回追加したい逆方向のために必要な補題は概ね次の二つに集約できる。
        \begin{bodyenumerate}
            \item \textbf{重みの一致を回収する補題}\\
            「$j\neq i$ の成分が一致し、かつ和が $1$ なら $i$ 成分も一致する」という事実を、
            \texttt{Set.indicator} と \texttt{Finset.sum} の言葉で整理する。
            \item \textbf{重みの一致 $\Rightarrow$ 直線所属}\\
            \texttt{Finset.affineCombination} が重みに関してアフィンであること（\texttt{lineMap} と可換）を用いて、
            「ある $t_i$ で展開した重み」から \texttt{AffineMap.lineMap} 形の点表示を得る。
        \end{bodyenumerate}

        \indent これらを組み合わせれば、定理\ref{thm:ceva-converse-target}は
        \texttt{Ceva.lean} と同程度の依存（\texttt{Simplex.Basic} まで）で実装できる見込みがある。
    \end{body}

    \section{Lean での定理名とスケッチ}\label{sec:lean-sketch}
    \begin{body}
        \indent 実際に追加する定理は、既存の命名に合わせて
        \begin{quote}
            \texttt{AffineIndependent.mem\_line\_of\_affineCombination\_eq\_smul\_eq}
        \end{quote}
        のような形が自然である（名前は最終的に mathlib コミュニティの慣習に従って調整する）。
        Lean の型としては概ね次のようになる。
    \end{body}

\begin{verbatim}
namespace AffineIndependent
variable [Ring k] [AddCommGroup V] [Module k V] [AffineSpace V P]
variable {ι : Type*} {p : ι → P}

-- (existing) forward direction
#check AffineIndependent.exists_affineCombination_eq_smul_eq

-- (new) converse direction (shape)
theorem mem_line_of_affineCombination_eq_smul_eq
    (hp : AffineIndependent k p) {s : Set ι} (hs : s.Nonempty)
    {fs : s → Finset ι} {w : s → ι → k} (hw : ∀ i, ∑ j ∈ fs i, w i j = 1)
    {p' : P} {w' : ι → k} {fs' : Finset ι}
    (hw' : (∑ j ∈ fs', w' j = 1) ∧ fs'.affineCombination k p w' = p')
    (hscale : ∀ i : s, ∃ r, ∀ j,
      r * Set.indicator ((fs i : Set ι) \ {(i : ι)}) (w i) j
        = Set.indicator ((fs' : Set ι) \ {(i : ι)}) w' j) :
    ∀ i : s, p' ∈ line[k, p i, (fs i).affineCombination k p (w i)] := by
  -- proof: reconstruct the lineMap parameter from hscale,
  -- show the induced weights agree with w' using hp (uniqueness),
  -- then conclude membership via mem_affineSpan_pair_iff_exists_lineMap_eq.
  ...
end AffineIndependent
\end{verbatim}

    \begin{body}
        \indent この形は「既存定理の結論をそのまま仮定に置く」だけなので、
        ユーザは順方向定理と組み合わせて必要な同値形を構成できる。

        \indent さらに古典的チェバの「平行を含む逆」を mathlib4 で美しく扱うには、
        射影化（無限遠点の導入）か、あるいは「重みの総和が $0$ の場合は直線が平行になる」という補題を追加して
        \emph{場合分け}を作る必要がある。
        本研究ではまずアフィン幾何の枠内で完結する上記の回収補題を優先し、
        平行を含む形への拡張は今後の課題として整理する。
    \end{body}


    % =========================================================
    \chapter{結論と今後の展望}
    % =========================================================

    \section{結論}
    \begin{body}
        \indent 既にJoseph Myers氏によって形式化された一般化チェバの定理を参考に、一般化チェバの定理の逆を形式化し、mathlib4にPRを提出した。
    \end{body}

    \section{今後の課題}
    \begin{body}
        \indent 当初は一般化チェバの定理の順方向を実装する予定だったが、コミュニケーション不足によりJoseph Myers氏に先行されることになった。研究内容の共有がいかに大切かを痛感した経験となった。

        \indent また、数学にしろ情報学にしろ、何を達成したいかを正確に把握して最適な言語・理論を選択することが不可欠であることを痛感した。単一の定理の証明を紙面上で構築するための概念の定義なのか理論の一部としてライブラリに用意する概念の定義なのかという違い、自然言語で定義するのか一階述語論理で定義するのか型理論で定義するのかという違い、によって最適な定義が揺れるためその選定に最も苦戦した。
        
        \indent 今回培ったOSS開発の作法に従い、次回以降のプロジェクトにも取り組んでいきたい。
    \end{body}

    \section{Lean4の普及と形式化研究の展望}
    \begin{body}
        \indent 現在LeanのAutoformalizationという技術開発が進んでおり、自然言語で記述した数学の証明をそのままLeanの形式言語に書き換えることが目指されている。\\
        今回は個別の定理を形式化すると言うプロジェクトであったが、この技術を用いればある数学的な証明を自然言語で記述するだけで、立ちどころに形式証明が得られると考えられている。
        この技術は、Lean研究を飛躍的に促進すると考えられる。個別の定理の形式化だけでなく、このAutoformalizationに関する見識も深めていきたい。
    \end{body}

    % =========================================================
    \appendix
    % =========================================================

    % =========================================================
    % 付録：型理論についての補足
    % =========================================================
    \chapter{型理論についての補足}\label{appendix:type-theory}

    \section{本付録の目的}
    \begin{body}
        \indent 本章で触れたように型理論では、命題を型・証明をその型の項、として表現する。しかしこれだけでは主に一階述語論理ベースの集合論に慣れている読者にはその全貌は見せられないように思う。よってよりイメージのつきやすい型理論の説明を一階述語論理と比較し、形式体系とは何かに立ち返りながら行うことが本付録の目的である。
    \end{body}

    \section{形式体系とは}

    \begin{body}
        \indent 数学の推論を厳密に形式化して分析する分野を数理論理学と呼び、形式体系とはその中で用いる概念の一つである。数理論理学は構文論と意味論に分かれるが、両者は形式体系を別のアプローチで分析する。構文論とは記号列の操作（導出）の規則に着目する手法であり、意味論とは記号列に解釈を与えてその真偽を論じる手法である。

        \indent 一階述語論理や高階述語論理、型理論など、形式体系と呼ばれるものは様々あり、形式言語を用いて数学の証明・推論を様々な方法で表現する。構文論的な観点で見たときこれらの多くに共通する眼目は、証明という対象を有限のデータとして扱う点にある、つまり\emph{証明とは有限の記号列でなければならない}。
    
        \indent 形式体系はしばしば \emph{言語}、\emph{公理}、\emph{推論規則}の三つ組として記述される。言語は「許される記号と、それらから整式（文法を満たす記号列）を作る規則」を与える。公理は「証明の出発点として採用する式またはその集合」である。推論規則は「前提となる式から結論となる式を得る規則」である。これら自体は有限である必要はない。実際、変数記号は自然数の個数だけ無数に用意できるし、公理はスキームという形で実質無限本存在する、その上推論規則の中には$\omega$規則$\frac{\{\varphi(n)\}_{n\in\mathbb{N}}}{\forall n\in\mathbb{N}\,\varphi(n)}\;(\omega)$と呼ばれるものもあり、推論が無限を含んでいる。

        \indent しかし証明を構築するときは話が別である。証明とは「人間が数学的主張の正しさを検証するための道具」であるから有限でなくては検証ができない。よっていざ証明を構築するときは、言語から有限個の記号を選び、公理スキームの有限個を利用し、有限個の仮定から有限個の結論が得られる推論規則のみを利用する。したがって形式体系を構築するにあたって、言語の無限個の記号、公理の無限本の公理図式は許容されるが、$\omega$規則のような\emph{無限推論規則}は「証明が有限」という前提自体を捨てる必要が生まれるため一般的な形式体系の議論では許容されない。（これを認めた無限論理・無限証明論と呼ばれる分野も存在するが、無限を扱える代わりに完全性やコンパクト性などの豊かな性質の多くが失われるまたは変質してしまう別物として区別され研究されている）
    
        \indent 本稿では、与えられた公理（集合）に推論規則を適用して得られる導出可能な式全体の閉包を \emph{理論体系}と呼ぶ。「これらの式が何を意味するか」という解釈を与えるのは意味論に譲る。
    
        \indent 例えば、多くの数学書が暗黙のうちに議論に用いている集合論（典型的には ZF / ZFC）は、形式体系としては「集合 $\in$ を含む言語」「公理（外延性・分出・置換スキーム等）」「古典論理の推論規則」を備えた理論として以下のように整理できる。
        
        \vskip\baselineskip

        \noindent\textbf{集合論}
        \begin{bodyitemize}[leftmargin=7em]
            \item[言\hspace*{2em}語] : 一階述語論理記号$\cup$変数記号$\cup$補助記号$\cup\in$と論理式の定義
            \item[推論規則] : 自然演繹
            \item[公\hspace*{2em}理] : ZF / ZFC
        \end{bodyitemize}

        \vskip\baselineskip

        \indent まず、言語を定義する。記号集合$\{\neg,\to,\forall,\exists$,…$\}\cup \{v_{0},v_{1},v_{2}$,…$\}\cup \{\in\}$を用意し、この集合から$\neg v_{2}\forall, \forall v_{0}(v_{0}=v_{0}), \cup v_{n}\cap v_{n-1}=v_{0}$のように無数の記号列を構成できる。そして得られた記号列から記号列として認めるものを定義する。認めた記号列を論理式と呼び$\phi,\psi,\varphi$などと略記する。先ほどの例では$\forall x(x=x)$のみ論理式である。\\
        次に、推論規則を定義する。有名な例では含意除去$\infer[\to E]{\varphi \\ \varphi \to \psi}{\psi}$（Modus Ponensのこと）などがある。先ほどの例に対しては例えば全称除去$\infer[\forall E]{\forall x\,\varphi(x)}{\varphi(t)}$を用いて$v_{0}=v_{0}$などの新たな論理式が得られる。\\
        最後に公理を決める。Kunenの定義に従えば、外延性の公理
        \[
        \forall v_0\,\forall v_1\;\Bigl(\forall v_2\,(v_2\in v_0 \leftrightarrow v_2\in v_1)\;\rightarrow\; v_0=v_1\Bigr)
        \]
        などがある。論理式の中で証明の出発点として採用する特別な論理式である。

        \indent こうして得られた公理から推論規則による変形を通して得られたものだけが定理と呼ばれ、理論体系を構成する一員となる。

        \indent こう書くと、自然に「では言語、公理、推論規則のいずれかを変更したらどうなるだろう？」という疑問が自然に湧いてくるであろう。公理的集合論では主にこの公理の部分を様々に入れ替え、付け加え、除くことで理論体系がどう変化するかを分析する。また、推論規則としてはヒルベルト式やシーケント計算が存在し定理の証明のしやすさに影響を与える。

        \indent 同様にしてこの三つ組を変更して得られるのが、型理論である。（正確にいうと、上記の三つ組は一階述語論理ベースの形式体系を分類するには有用だが、型理論ベースの形式体系を分類するには不向きである。その上で敢えて同じ視点でCICを眺めてみる。）

        
        \vskip\baselineskip

        \noindent\textbf{CIC（Calculus of Inductive Constructions）を同じ視点で眺める}\par
        \vskip0.5\baselineskip

        \indent Lean（および Coq/Rocq）の基礎にある体系は、CIC と呼ばれる依存型理論の系統である。
        ただし、集合論のように「言語＋公理＋推論規則」という三つ組で捉えるよりも、
        \emph{判断（judgement）とその生成規則}として捉える方が理解しやすい。
        ここでは集合論に慣れた読者が迷いやすい点に絞って概観する。

        \begin{bodyitemize}[leftmargin=4.5em]
            \item[\textbf{（1）判断が基本単位}]\
            一階論理で「$\vdash \varphi$」の形を扱うのに対し、CIC では典型的に
            \[
              \Gamma\ \mathsf{ctx},\qquad
              \Gamma \vdash A\ \mathsf{type},\qquad
              \Gamma \vdash t : A
            \]
            のような判断を生成する規則を与える。$\Gamma$ は文脈（変数宣言の列）であり、
            一階論理の「仮定集合」に似ているが、\emph{文脈自体も推論の対象}である点が異なる。

            \item[\textbf{（2）$\Pi$ 型は「全称」と「含意」を統合する}]\
            依存積（$\Pi$ 型）$\Pi x:A.\,B(x)$ は、$x$ に依存する型 $B(x)$ を許す関数型である。
            依存しない場合 $\Pi x:A.\,B$ は $A\to B$ と同じで、自然演繹における
            「仮定導入／仮定消去」と対応する（Curry--Howard対応）。
            ここが「命題＝型、証明＝項」という読み方の出発点になる。

            \item[\textbf{（3）等号が二種類ある}]\
            Lean のカーネルが自動で判定する等しさは、主に\emph{定義的等しさ}（definitional equality）である。
            これは $\beta$ 簡約（$(\lambda x.\,t)\,u$ の計算）や帰納型の計算規則で同一視してよいという約束で、
            型検査に組み込まれている。
            一方、\emph{命題としての等号}は \texttt{Eq}（同一視型）として表され、その証明は通常の証明と同様に構成が必要である。
            「計算で勝手に等しくなる」部分と「証明が必要な等号」を区別することが、集合論出身者にとって重要な注意点である。

            \item[\textbf{（4）帰納型は「公理」ではなく「生成規則」として入る}]\
            自然数 $\mathbb N$ やリスト等は、生成子（\texttt{zero}, \texttt{succ} 等）を持つ\emph{帰納型}として導入され、
            その上で再帰・帰納法（除去規則）が与えられる。
            集合論で言えば「無限公理」などを通して $\omega$ を確保し、その上で再帰を展開するのに対し、
            型理論では「この型はこう生成され、こう消去できる」という規則が体系に組み込まれる。

            \item[\textbf{（5）宇宙（universes）}]\
            「すべての型の型」を一つにまとめるとラッセル型の破綻が起きるため、
            CIC では $Type_0 : Type_1 : Type_2 : \cdots$ の階層（宇宙）を用いる。
            これは集合論の累積階層に似た安全装置だが、対象は集合ではなく型である。
            Lean の \texttt{Type*} は「ある宇宙 \texttt{Type u} のどれか」を意味する略記である。
        \end{bodyitemize}

        \vskip\baselineskip

        \indent 以上を踏まえると、CIC は「ある対象世界（集合の宇宙）をまず仮定し、その上で真偽を語る」体系というより、
        「何が項として許され、どの型を持つかを決める\emph{生成規則の体系}」であると言える。
        その上で、命題を型として読み、証明を項として構成することで、数学の議論が Lean のコードとして記述される。

        \indent なお、CIC でも必要に応じて \emph{公理を追加}できる。
        Lean では「ある型 $A$ に住む定数 $c:A$ を宣言する」ことが公理追加に相当する。
        mathlib4 では選択公理や古典論理（排中律）を利用する場面も多く、
        その意味で \emph{型理論＝必ず構成的} というわけではない。
        ただし、どの公理を足すかは計算可能性や抽出可能性に影響するため、
        形式化の設計としては意識的に管理する必要がある。

        \indent 型理論は一階述語論理とは異なる形で「証明」を有限データとして取り扱い、
        カーネルによる機械検証を可能にする。これが、Lean4/mathlib4 で数学を形式化する際の基盤である。

    \end{body}

    % =========================================================
    % 付録：　古典的チェバの定理の版の教科書による違いについて
    % =========================================================
    \chapter{古典的チェバの定理の版の教科書による違いについて}\label{appendix:classical-ceva-version}

    \section{本付録の目的}
    \begin{body}
        \indent 本付録では、高校数学で学ぶ古典的チェバの定理の版の違いがなぜ生まれ、どのような意図があるのかを解説する。
    \end{body}

    \section{古典的チェバの定理の版の種類}
    \begin{body}
        \indent 古典的チェバの定理は大きく分けて以下の3つの形で紹介されうる。
        \begin{bodyitemize}
            \item[①] $D, E, F$を辺上に限定した、3直線が1点で交わるための必要十分条件
            \item[②] $D, E, F$を辺の延長上でも許した、3直線が1点で交わる必要条件(定理\ref{thm:classical-ceva})
            \item[③] $D, E, F$を辺の延長上でも許した、3直線が1点で交わる必要十分条件
        \end{bodyitemize}

        \indent 一見それぞれ異なる主張をしているように見えるがこれらは数学的に何も問題はなく、\emph{背後の幾何学体系に何を採用するか}で主張の範囲や述べ方が違うだけである。

    \end{body}

    \section{古典的チェバに版の違いが生まれる要因}
    \begin{body}
        \indent ①〜③の中で本稿で取り扱っている②のみ必要十分性を主張する定理になっていない。試しに定理の安直な逆命題を考えてみると成立しないことがわかる。その原因となっているのが、下図のような、チェビアンが平行になる場合である。
    \end{body}

    \begin{center}
        \begin{tikzpicture}[scale=0.8, x=1cm,y=1cm, line cap=round, line join=round, >=Latex]
            % ---- keep the whole picture nicely centered ----
            \path[use as bounding box] (-1.2,-1) rectangle (7.2,10.2);
    
            % --- Title (Added per request) ---
            \node[font=\bfseries] at (3,9.8) {(図)等式条件を満たしながらチェビアンが平行となる場合};
    
            % --- coordinates ---
            \coordinate (A) at (0,0);
            \coordinate (B) at (6,0);
            \coordinate (C) at (4,3);
    
            \coordinate (D) at (0,9);
            \coordinate (E) at (6,4.5);
            \coordinate (F) at (4,0);
    
            % --- triangle ---
            \draw[thick] (A) -- (B) -- (C) -- cycle;
    
            % --- side extensions (dashed) ---
            \draw[dashed] (B) -- (D);  % extension of BC to meet x=0
            \draw[dashed] (A) -- (E);  % extension of CA to meet x=6
    
            % --- points ---
            \fill (A) circle (1.6pt) node[below left] {$A$};
            \fill (B) circle (1.6pt) node[below right] {$B$};
            \fill (C) circle (1.6pt) node[above, xshift=2pt] {$C$};
    
            \fill (D) circle (1.6pt) node[left] {$D$};
            \fill (E) circle (1.6pt) node[right] {$E$};
            \fill (F) circle (1.6pt) node[below] {$F$};
    
            % --- Improved Parallel Mark Style ---
            \tikzset{
              parallelArrow/.style={
                postaction={decorate},
                decoration={
                  markings,
                  mark=at position 0.5 with {
                    \draw[line width=0.8pt] (-3pt,-3pt) -- (0pt,0pt) -- (-3pt,3pt);
                  }
                }
              }
            }
    
            % --- cevians (parallel) with fixed arrows ---
            \draw[very thick, parallelArrow] (A) -- (D);
            \draw[very thick, parallelArrow] (B) -- (E);
            \draw[very thick, parallelArrow] (F) -- (C); 
    
        \end{tikzpicture}
    \end{center}

    \begin{body}
        \indent この場合、等式条件は成立するがチェビアンが平行となってしまい共点していない。まさにこの問題によって、背景とする幾何学体系による定理の述べ方が変わるのである。
    \end{body}

    \section{各々の版の背景}

    \begin{body}
        \indent 結論から述べると、①〜③はそれぞれユークリッド幾何、アフィン幾何、射影幾何の観点から古典的チェバの定理を述べたものである。

        \indent ①はユークリッド幾何の観点で述べた形である、$BD$、$DC$などは全てユークリッド距離であり常に正となる。ユークリッド距離ではチェビアンの足が対辺を外分するときをうまく表現できないため、あらかじめチェビアンの足は対辺中に限定される。チェビアンが平行になることは排除できないが、その場合等式条件が成立しないため必要十分性を主張する定理となる。

        \indent ②はアフィン幾何の観点で述べた形である。ここで登場するのはアフィン比と呼ばれる$\dfrac{BD}{DC}$という値であり、辺を成す$2$頂点の重心座標から得られるものである。高校範囲内ではアフィン幾何を登場させられないため、ユークリッド幾何の範疇で定義できる有向線分という考え方を用いて表現している。しかしこの場合定理の逆がそのままでは成り立たず、等式条件が成立してもチェビアンが平行になってしまうという反例が生じる。

        \indent ③は射影幾何の観点で述べた形である。射影幾何はユークリッド幾何やアフィン幾何を包含する幾何学の大きな枠組みの一つである。重要な特徴が空間内に無限遠点を含むため、平行という概念を「無限遠点で交点を持つ」と言える点である。定理の主張自体は②と変わっていないが、交点をもつという言葉の意味が広がったことで述べ方が変わっている。

        \indent 高校数学で習うのは高々ユークリッド幾何であり違いを説明するのが困難なため、追加の前提条件を加えるなどして主張が偽にならないように工夫して解説するという教育的な配慮がなされている。

        \indent 本稿ではアフィン幾何に基づいて議論を進めた。「チェビアンが平行でない」という趣旨の前提を加えれば強引に必要十分性を主張する定理にできるが、定理の自然さが損なわれると考え②の形を選択した。
    \end{body}

    \backmatter
    \addchaptertotoc{謝辞}
    \begin{body}
        % TODO
    \end{body}

    \begin{thebibliography}{99}
        \bibitem{Kunen}Kenneth Kunen, \emph{Set Theory: An Introduction to Independence Proofs}, North-Holland (Elsevier Science Publishers B.V.), 1980.
        \bibitem{LeanProverCommunity} Lean Prover Community, \emph{Lean Prover Community}, \url{https://leanprover-community.github.io/}, accessed 2026-01-18.
        \bibitem{DeepMindIMO2025} Google DeepMind, \emph{Advanced version of Gemini with Deep Think officially achieves gold-medal standard at the International Mathematical Olympiad}, 2025.
        \bibitem{multipede} Dov Samet, \emph{An Extension of Ceva's Theorem ton-Simplices}, 2021.
        \bibitem{LeanResearch} 秋田 隼, \emph{LeanResearch}, GitHub repository, \url{https://github.com/Aj1905/LeanResearch.git}, 2026.
        \bibitem{mathlib3_ceva} Mantas Bakšys, \emph{Ceva's Theorem (mathlib3)}, GitHub Pull Request \#10632, \url{https://github.com/leanprover-community/mathlib3/pull/10632}, 2021--2023.
        \bibitem{mathlib4_ceva_alexeev} Boris Alexeev, \emph{Ceva's Theorem (mathlib4)}, GitHub Pull Request \#33388, \url{https://github.com/leanprover-community/mathlib4/pull/33388}, 2025--2026.
        \bibitem{mathlib4_ceva_myers} Joseph Myers, \emph{Generalized Ceva's Theorem in $n$-dimensional affine space (mathlib4)}, GitHub Pull Request \#33409, \url{https://github.com/leanprover-community/mathlib4/pull/33409}, 2026.

        \vskip\baselineskip

        また、研究にあたりOpenAI社のChat GPT 5.2 thinkingモデルを大いに活用した。以下は使用したチャット履歴の一部である。\\
        \url{https://chatgpt.com/g/g-p-696f0c41aa3c8191928eeed4b3fe9bb8-mathlib4-affinespace/project}
    \end{thebibliography}

\end{document}
